\documentclass[a4paper,openright,12pt]{report}
\usepackage[spanish]{babel}			% Permite que partes automáticas del documento aparezcan en castellano.
\usepackage[utf8]{inputenc}			% Permite escribir tildes y otros caracteres directamente en el .tex
\usepackage[T1]{fontenc}   		         % Asegura que el documento resultante use caracteres de una fuente
\usepackage{listings}				% Permite utilizar lenguajes de programacion dentro de latex
\usepackage{enumerate}				%enumerados
\usepackage{graphicx}%graficos

\begin{document}

\begin{titlepage}

\begin{center}
\vspace*{-1in}
\begin{figure}[htb]
\begin{center}
\includegraphics[width=8cm]{./imagenes/espol.jpg}
\end{center}
\end{figure}
FACULTAD DE INGENIERIA EN ELECTRICIDAD Y COMPUTACIÓN\\
\vspace*{0.15in}
LENGUAJES DE PROGRAMACIÓN\\
\vspace*{0.6in}
\begin{large}
PROYECTO DE SEGUNDO PARCIAL: HASKELL
\end{large}
\vspace*{0.4in}
\begin{large}
ANALIZADOR SINTÁCTICO DE UN ARCHIVO XML\\
\end{large}
\vspace*{0.3in}
\rule{80mm}{0.1mm}\\
\vspace*{0.1in}
\begin{large}
Integrantes:\\Leonel Ramírez Gonzalez\\José Vélez Gómez\\Kevin Campuzano Castillo\\ 
\end{large}
\end{center}
\end{titlepage}

\tableofcontents
\chapter{Objetivos}
\begin{itemize}
\item {Aprender como funciona un lenguaje puramente funcional como es Haskell.}
\item{Recononcer como trabaja un analizador sintáctico, que es una de las partes de un compilador.}
\end{itemize}

\chapter{Introducción}
\textbf{Haskell} es un lenguaje de programación puramente funcional , no estricto y fuertemente tipado que fue disenado por la universidades de Yale y la universidades de Glasgow.\\

Nace en como la solución de la crisis de los años sesenta en la que la mayoria de software que se producía no era fiable, tenian una gran tasa de errores que ponian en grave peligro la confianza de los usuarios en estos sistemas por esta razón se creo este lenguaje como un nuevo modelo de programación al que se lo conoce como programación funcional.\\


\chapter{Alcance del Proyecto}


\chapter{Observaciones}
\begin{itemize}
\item \textbf{Leonel Ramírez Gonzalez}
\begin{itemize}
	\item[Ventajas: ]Una de las ventajas que me agrado en haskell fue que al momento de leer el archivo su implementación era mas sencilla de la que se implementa en otros lenguajes.
	\item[Desventajas: ]Una de las desventajas que encontre yo de Haskell fue la manera de como se reciben los parametros en las funciones.
\end{itemize}
\end{itemize}
\end{document}